\section*{Введение}
\addcontentsline{toc}{section}{Введение}

В современных маркетинговых исследованиях и анализе данных ключевым этапом является
сбор репрезентативной информации о предпочтениях потребителей. В условиях
ограниченности ресурсов проведение сплошного опроса генеральной совокупности
зачастую невозможно, что обуславливает актуальность применения различных
методоввыборочного наблюдения.

\textbf{Актуальность работы} заключается в необходимости поиска наиболее точных и
эффективных алгоритмов формирования выборки, которые позволяют минимизировать статистическую
погрешность при оценке потребительских качеств товаров и брендов.

\textbf{Объектом исследования} является массив данных опроса респондентов, содержащий
оценки характеристик ведущих аудио-брендов (Sony, JBL, Marshall) по различным критериям (качество звука, цена, дизайн и др.).

\textbf{Предметом исследования} выступают методы статистического отбора и их влияние
на точность оценки лояльности потребителей в рамках модели идеальной точки.

\textbf{Целью данной работы} является проведение сравнительного анализа алгоритмов
статистического отбора и обоснование выбора оптимальной методики формирования выборки,
обеспечивающей максимальную репрезентативность данных и устойчивость интегральных
показателей потребительской лояльности.

Для достижения поставленной цели необходимо решить следующие \textbf{задачи}:
\begin{enumerate}
    \item Осуществить предварительную \textbf{обработку данных} генеральной совокупности 
    респондентов.
    \item Реализовать алгоритмически $4$ метода отбора данных: \emph{случайный},
    \emph{механический}, \emph{типический} и \emph{кластерный}.
    \item Вычислить \textbf{интегральный показатель брендов}, используя модель 
    \emph{идеальной точки}, учитывающую веса параметров и ожидания потребителей.
    \item Провести статистическую оценку \textbf{репрезентативности} полученных выборок.
    \item Сравнить \textbf{эффективность методов} отбора на основе полученных
    показателей точности и визуализировать результаты анализа.
\end{enumerate}

Практическая значимость работы заключается в разработке инструментария 
на языке \texttt{Python}, который автоматизирует расчет рейтингов брендов 
и их погрешностей. Созданный программный код наглядно демонстрирует
разные методы отбора данных и позволяет выбрать тот, который дает наиболее точный результат
с минимальной ошибкой.